\section{Related Work}

\subsection{RLHF Evaluation and Benchmarks}

Reinforcement Learning from Human Feedback has become the standard approach for aligning language models with human preferences. \citet{ouyang2022training} introduced InstructGPT using RLHF to improve helpfulness and safety. Evaluation typically relies on aggregate benchmark metrics---TruthfulQA measures factual accuracy \citep{lin2022truthfulqa}, ETHICS assesses moral reasoning \citep{hendrycks2021ethics}, and HHH evaluates helpfulness, harmlessness, and honesty \citep{askell2021general}.

However, these benchmarks report aggregate accuracy without task-level analysis. Our work differs by clustering tasks by behavioral signals (calibration patterns), enabling mechanism investigation at the task-type level rather than aggregate statistics.

\subsection{Model Calibration}

Calibration measures alignment between model confidence and actual correctness \citep{guo2017calibration}. Prior work treats miscalibration as a general phenomenon. We show calibration inversion clusters systematically (silhouette = 0.6016), indicating specific task types trigger miscalibration rather than uniform overconfidence.

\subsection{Bidirectional Alignment Framework}

\citet{shen2024bidirectional} proposed a theoretical framework distinguishing AI-to-Human alignment from Human-to-AI alignment. Our work operationalizes this framework, hypothesizing that tasks requiring bidirectional adaptation correlate with calibration inversion. While our keyword-based feature detection proved insufficient, we verify the underlying mechanism chain.

\subsection{Positioning}

Prior work either evaluates RLHF outcomes or analyzes mechanisms in isolation. We bridge these by using behavioral patterns to investigate training mechanisms.
